\section{Perceptron wielowarstwowy}
\label{sec:mlp}

\subsection{Model sieci}
\label{sec:mlp-model}
% Plan (etap 4): notacja macierzowa, warstwa Z = A W + b (wiersz = próbka), bias
% opcjonalny, propagacja w przód.

\subsection{Funkcje aktywacji}
\label{sec:mlp-activations}
% Plan (etap 4): tabela sigmoid, tanh, ReLU, leaky ReLU (wzór, pochodna, zakres);
% wyjście: identity (regresja), softmax (klasyfikacja).

\subsection{Funkcje straty}
\label{sec:mlp-losses}
% Plan (etap 4): entropia krzyżowa i MSE na kodowaniu one-hot (klasyfikacja),
% MSE i MAE (regresja).

\subsection{Propagacja wsteczna}
\label{sec:mlp-backprop}
% Plan (etap 4): wzory gradientów warstwa po warstwie, skrót (P - Y)/n dla entropii
% krzyżowej z softmax, ogólny VJP softmaxu; algorytm.

\subsection{Optymalizacja}
\label{sec:mlp-optimization}
% Plan (etap 4): mini-batch SGD, momentum, Adam; tasowanie co epokę.

\subsection{Inicjalizacja wag}
\label{sec:mlp-init}
% Plan (etap 4): inicjalizacja Xaviera (sigmoid, tanh) i He (ReLU, leaky ReLU).
